% Options for packages loaded elsewhere
\PassOptionsToPackage{unicode}{hyperref}
\PassOptionsToPackage{hyphens}{url}
%
\documentclass[
]{article}
\usepackage{amsmath,amssymb}
\usepackage{iftex}
\ifPDFTeX
  \usepackage[T1]{fontenc}
  \usepackage[utf8]{inputenc}
  \usepackage{textcomp} % provide euro and other symbols
\else % if luatex or xetex
  \usepackage{unicode-math} % this also loads fontspec
  \defaultfontfeatures{Scale=MatchLowercase}
  \defaultfontfeatures[\rmfamily]{Ligatures=TeX,Scale=1}
\fi
\usepackage{lmodern}
\ifPDFTeX\else
  % xetex/luatex font selection
\fi
% Use upquote if available, for straight quotes in verbatim environments
\IfFileExists{upquote.sty}{\usepackage{upquote}}{}
\IfFileExists{microtype.sty}{% use microtype if available
  \usepackage[]{microtype}
  \UseMicrotypeSet[protrusion]{basicmath} % disable protrusion for tt fonts
}{}
\makeatletter
\@ifundefined{KOMAClassName}{% if non-KOMA class
  \IfFileExists{parskip.sty}{%
    \usepackage{parskip}
  }{% else
    \setlength{\parindent}{0pt}
    \setlength{\parskip}{6pt plus 2pt minus 1pt}}
}{% if KOMA class
  \KOMAoptions{parskip=half}}
\makeatother
\usepackage{xcolor}
\usepackage[margin=1in]{geometry}
\usepackage{graphicx}
\makeatletter
\def\maxwidth{\ifdim\Gin@nat@width>\linewidth\linewidth\else\Gin@nat@width\fi}
\def\maxheight{\ifdim\Gin@nat@height>\textheight\textheight\else\Gin@nat@height\fi}
\makeatother
% Scale images if necessary, so that they will not overflow the page
% margins by default, and it is still possible to overwrite the defaults
% using explicit options in \includegraphics[width, height, ...]{}
\setkeys{Gin}{width=\maxwidth,height=\maxheight,keepaspectratio}
% Set default figure placement to htbp
\makeatletter
\def\fps@figure{htbp}
\makeatother
\setlength{\emergencystretch}{3em} % prevent overfull lines
\providecommand{\tightlist}{%
  \setlength{\itemsep}{0pt}\setlength{\parskip}{0pt}}
\setcounter{secnumdepth}{-\maxdimen} % remove section numbering
% definitions for citeproc citations
\NewDocumentCommand\citeproctext{}{}
\NewDocumentCommand\citeproc{mm}{%
  \begingroup\def\citeproctext{#2}\cite{#1}\endgroup}
\makeatletter
 % allow citations to break across lines
 \let\@cite@ofmt\@firstofone
 % avoid brackets around text for \cite:
 \def\@biblabel#1{}
 \def\@cite#1#2{{#1\if@tempswa , #2\fi}}
\makeatother
\newlength{\cslhangindent}
\setlength{\cslhangindent}{1.5em}
\newlength{\csllabelwidth}
\setlength{\csllabelwidth}{3em}
\newenvironment{CSLReferences}[2] % #1 hanging-indent, #2 entry-spacing
 {\begin{list}{}{%
  \setlength{\itemindent}{0pt}
  \setlength{\leftmargin}{0pt}
  \setlength{\parsep}{0pt}
  % turn on hanging indent if param 1 is 1
  \ifodd #1
   \setlength{\leftmargin}{\cslhangindent}
   \setlength{\itemindent}{-1\cslhangindent}
  \fi
  % set entry spacing
  \setlength{\itemsep}{#2\baselineskip}}}
 {\end{list}}
\usepackage{calc}
\newcommand{\CSLBlock}[1]{\hfill\break\parbox[t]{\linewidth}{\strut\ignorespaces#1\strut}}
\newcommand{\CSLLeftMargin}[1]{\parbox[t]{\csllabelwidth}{\strut#1\strut}}
\newcommand{\CSLRightInline}[1]{\parbox[t]{\linewidth - \csllabelwidth}{\strut#1\strut}}
\newcommand{\CSLIndent}[1]{\hspace{\cslhangindent}#1}
\usepackage{float}
\ifLuaTeX
  \usepackage{selnolig}  % disable illegal ligatures
\fi
\usepackage{bookmark}
\IfFileExists{xurl.sty}{\usepackage{xurl}}{} % add URL line breaks if available
\urlstyle{same}
\hypersetup{
  pdftitle={Homework 5},
  hidelinks,
  pdfcreator={LaTeX via pandoc}}

\title{Homework 5}
\author{}
\date{\vspace{-2.5em}}

\begin{document}
\maketitle

\subsubsection{Definitions}\label{definitions}

Define the following terms/concepts

\begin{itemize}
\item
  a.) \textbf{Non-parametric statistics}
\item
  b.) \textbf{Parametric statistics}
\item
  c.) Describe the advantages and disadvantages of
  \textbf{Non-parametric} significance tests
\item
  d.) \textbf{Explanatory Variable}
\item
  e.) \textbf{Response Variable}
\end{itemize}

\subsubsection{Confidence intervals and two-tailed
tests:}\label{confidence-intervals-and-two-tailed-tests}

\(\bf (1.)\) For parts (a) - (d) you are given a confidence interval for
the parameter along with the point estimate it was built from. Determine
whether the null hypothesis \(H_0\) should be \textbf{Rejected} or
\textbf{Not Rejected} by checking whether the value claimed under
\(H_0\) falls inside the interval, and give the significance level
\(\alpha\)

a.) \(H_0: \mu = -2\) \(H_A: \mu \neq -2\) \(95\%\) CI for
\(\mu \in [-3.0, -2.4]\) \(\bar{x} = -2.7\)
\[ \text{Decision} = \_\_\_\_\_\_\_ \\
  \alpha = \_\_\_\_\_\_ \]

b.) \(H_0: p = 0.9\) \(H_A: p \neq 0.9\) \(99\%\) CI for
\(p \in [0.968, 0.972]\) \(\hat{p} = 0.97\)

\[ \text{Decision} = \_\_\_\_\_\_\_ \\
  \alpha = \_\_\_\_\_\_ \]

c.) \(H_0: \mu = 0\) \(H_A: \mu \neq 0\) \(90\%\) CI for
\(\mu \in [-3.63, 7.49]\) \(\bar{x} = 1.93\)

\[ \text{Decision} = \_\_\_\_\_\_\_ \\
  \alpha = \_\_\_\_\_\_ \]

d.) \(H_0: \mu = 50\) \(H_A: \mu \neq 50\) \(85\%\) CI for
\(\mu \in [93.6, 139]\) \(\bar{x} = 116.3\)

\[ \text{Decision} = \_\_\_\_\_\_\_ \\
  \alpha = \_\_\_\_\_\_ \]

\subsubsection{The sign test}\label{the-sign-test}

\(\bf (2.)\) A dermatologist is conducting an experiment to test whether
a new topical skin care cream reduces the appearance of facial acne.
From a sample of \(20\) patients with acne the researcher gives each
patient a score from \(1-5\), where (1) represents the mildest acne and
(5) represents the most severe acne, before the topical cream is
applied. Patients are then asked to use the topical cream for \(10\)
days. Following the \(10\) day treatment, patients are given a new acne
score between \(1-5\). The data for this experiment is given in the
table below:

\begin{table}[H]
\centering\centering
\begin{tabular}{rrrll}
\toprule
Patient Number & Score Before Treatment & Score After Treatment & Difference & Sign\\
\midrule
\cellcolor{gray!10}{1} & \cellcolor{gray!10}{3} & \cellcolor{gray!10}{2} & \cellcolor{gray!10}{} & \cellcolor{gray!10}{}\\
2 & 2 & 1 &  & \\
\cellcolor{gray!10}{3} & \cellcolor{gray!10}{4} & \cellcolor{gray!10}{1} & \cellcolor{gray!10}{} & \cellcolor{gray!10}{}\\
4 & 4 & 1 &  & \\
\cellcolor{gray!10}{5} & \cellcolor{gray!10}{4} & \cellcolor{gray!10}{1} & \cellcolor{gray!10}{} & \cellcolor{gray!10}{}\\
\addlinespace
6 & 5 & 2 &  & \\
\cellcolor{gray!10}{7} & \cellcolor{gray!10}{1} & \cellcolor{gray!10}{1} & \cellcolor{gray!10}{} & \cellcolor{gray!10}{}\\
8 & 5 & 4 &  & \\
\cellcolor{gray!10}{9} & \cellcolor{gray!10}{5} & \cellcolor{gray!10}{1} & \cellcolor{gray!10}{} & \cellcolor{gray!10}{}\\
10 & 2 & 1 &  & \\
\addlinespace
\cellcolor{gray!10}{11} & \cellcolor{gray!10}{2} & \cellcolor{gray!10}{2} & \cellcolor{gray!10}{} & \cellcolor{gray!10}{}\\
12 & 5 & 1 &  & \\
\cellcolor{gray!10}{13} & \cellcolor{gray!10}{3} & \cellcolor{gray!10}{3} & \cellcolor{gray!10}{} & \cellcolor{gray!10}{}\\
14 & 3 & 1 &  & \\
\cellcolor{gray!10}{15} & \cellcolor{gray!10}{1} & \cellcolor{gray!10}{1} & \cellcolor{gray!10}{} & \cellcolor{gray!10}{}\\
\addlinespace
16 & 2 & 1 &  & \\
\cellcolor{gray!10}{17} & \cellcolor{gray!10}{4} & \cellcolor{gray!10}{1} & \cellcolor{gray!10}{} & \cellcolor{gray!10}{}\\
18 & 5 & 1 &  & \\
\cellcolor{gray!10}{19} & \cellcolor{gray!10}{2} & \cellcolor{gray!10}{3} & \cellcolor{gray!10}{} & \cellcolor{gray!10}{}\\
20 & 1 & 1 &  & \\
\bottomrule
\end{tabular}
\end{table}

The dermatologist is particularly interested in whether the new product
improves the appearance of the patients' acne. They plan to test the
null hypothesis \(H_0: p = 0.5\) using a right-tailed test with
alternative hypothesis \(H_A: p>0.5\) using a sign test to determine if
there is enough evidence to conclude that the topical cream has a
positive effect on reducing facial acne at the \(\alpha = 0.05\) level.

a.) Fill in the table above by computing the difference in acne scores
before and after applying the skin care cream. Record the sign of each
difference and report the total number of positive signs \(s\). Note
that any patient whose score did not change (a difference of zero) is
dropped from the test, so the number of observations used in the sign
test is the number of non-zero differences.

b.) What does a ``positive'' sign indicate in this experiment?

c.) Compute the \(p\)-value for the sign test

d.) Use the \(p\)-value from part (c) to determine if there is enough
evidence to conclude that the topical cream product has a positive
effect on reducing the appearance of acne. Interpret your decision in
context.

\(\bf (3.)\) In a boreal forest, researchers conducted an experiment to
study how herbivores respond to variations in food plant quality and
quantity at the stand level. They fertilized young forest stands and
observed herbivore use over the subsequent year. The data collected
included the number of animal tracks in the fertilized and control plots
(Ball, Danell, and Sunesson 2000)

\begin{table}[H]
\centering\centering
\begin{tabular}{rrrll}
\toprule
Observation & Number of Tracks In Fertilized plots & Number of Tracks In Control Plots & Difference & Sign\\
\midrule
\cellcolor{gray!10}{1} & \cellcolor{gray!10}{15} & \cellcolor{gray!10}{10} & \cellcolor{gray!10}{} & \cellcolor{gray!10}{}\\
2 & 12 & 9 &  & \\
\cellcolor{gray!10}{3} & \cellcolor{gray!10}{18} & \cellcolor{gray!10}{11} & \cellcolor{gray!10}{} & \cellcolor{gray!10}{}\\
4 & 14 & 8 &  & \\
\cellcolor{gray!10}{5} & \cellcolor{gray!10}{16} & \cellcolor{gray!10}{12} & \cellcolor{gray!10}{} & \cellcolor{gray!10}{}\\
\addlinespace
6 & 13 & 10 &  & \\
\cellcolor{gray!10}{7} & \cellcolor{gray!10}{17} & \cellcolor{gray!10}{11} & \cellcolor{gray!10}{} & \cellcolor{gray!10}{}\\
8 & 14 & 9 &  & \\
\cellcolor{gray!10}{9} & \cellcolor{gray!10}{19} & \cellcolor{gray!10}{13} & \cellcolor{gray!10}{} & \cellcolor{gray!10}{}\\
10 & 15 & 10 &  & \\
\addlinespace
\cellcolor{gray!10}{11} & \cellcolor{gray!10}{11} & \cellcolor{gray!10}{8} & \cellcolor{gray!10}{} & \cellcolor{gray!10}{}\\
12 & 16 & 12 &  & \\
\bottomrule
\end{tabular}
\end{table}

The researchers are interested in whether or not fertilizing stands
increased herbivore activity. Conduct a sign test at the
\(\alpha = 0.05\) significance level to determine if fertilized plots
have significantly more herbivore tracks than control plots

a.) Fill in the table above by computing the difference in animal tracks
between the fertilized and control plots. Record the sign of each
difference and report the total number of positive signs \(s\)

b.) What does a ``positive'' sign indicate in this experiment?

c.) Compute the \(p\)-value for the sign test

d.) Use the \(p\)-value from part (c) to determine if there is enough
evidence to conclude that fertilized stands have more herbivore
activity.

\subsubsection{\texorpdfstring{Comparing two population proportions
\(p_1\) and
\(p_2\)}{Comparing two population proportions p\_1 and p\_2}}\label{comparing-two-population-proportions-p_1-and-p_2}

\(\bf (4.)\) The United States Supreme Court serves as the judicial
branch of the U.S. government, tasked with ensuring that laws conform to
the U.S. Constitution. Although traditionally operating as a relatively
discreet arbiter within the American government, recent years have seen
an unusual spotlight on the Supreme Court in partisan politics. This
heightened attention is largely attributed by political analysts to
several landmark Supreme Court decisions and ethical controversies that
have surfaced. Consequently, it is believed that this increased scrutiny
has adversely affected public approval of the Court across the political
spectrum. The table below provides a summary of current and historical
polling results for assessing SCOTUS job approval. The current approval
rating is summarized from \(192\) political polls conducted among the
public since late 2022 (Ryan and Bycoffe 2024). The historic approval
sentiment is characterized from \(39\) Gallup polls conducted since
early 2001.

\begin{table}[H]
\centering\centering
\begin{tabular}{llrr}
\toprule
Period & Parameter & Sample Size & Mean SCOTUS Approval Rating\\
\midrule
\cellcolor{gray!10}{Current Approval (2022 - 2024)} & \cellcolor{gray!10}{$p_1$} & \cellcolor{gray!10}{192} & \cellcolor{gray!10}{41.16}\\
Past Approval (Gallup: 2001 - 2022) & $p_2$ & 39 & 50.87\\
\bottomrule
\end{tabular}
\end{table}

Conduct a two-sample test at the \(\alpha = 0.05\) significance level to
determine whether the current approval rating is significantly lower
than the historical average - be sure to address all five steps of the
hypothesis test.

\(\bf (5.)\) A study published in the Journal of Paediatrics and Child
Health was interested in the number of reports of child abuse in
Sergipe, Brazil, before and after COVID-19 started. The researchers
hypothesized that children in abusive homes might be in more danger
because they weren't going to school and their families were under more
stress (Martins-Filho et al. 2020). A summary of the study is reported
below:

\begin{table}[H]
\centering\centering
\begin{tabular}{llrrr}
\toprule
Period & Parameter & Total Registries & Registries for Children under 12 Years of Age & Relative Proportion\\
\midrule
\cellcolor{gray!10}{Before COVID-19 (2019)} & \cellcolor{gray!10}{$p_1$} & \cellcolor{gray!10}{70} & \cellcolor{gray!10}{18} & \cellcolor{gray!10}{0.26}\\
After COVID-19 (2020) & $p_2$ & 53 & 16 & 0.30\\
\bottomrule
\end{tabular}
\end{table}

Conduct a two-sample test at the \(\alpha = 0.01\) significance level to
determine whether the proportion of abuse registries involving children
under the age of \(12\) was significantly different after the start of
the COVID-19 pandemic - be sure to address all five steps of the
hypothesis test. (Note that this proportion is the share of all
registries that involve a child under 12, not the underlying rate of
abuse in the population.)

\subsubsection{\texorpdfstring{Comparing two population means \(\mu_1\)
and
\(\mu_2\)}{Comparing two population means \textbackslash mu\_1 and \textbackslash mu\_2}}\label{comparing-two-population-means-mu_1-and-mu_2}

\(\bf (6.)\) A study published in the journal Physical Culture and
Sport. Studies and Research explored what motivates athletes to keep
playing sports, comparing team sports like football with individual ones
like taekwondo (Moradi, Bahrami, and Dana 2020). The study used
stratified random sampling to survey \(265\) athletes from four team
disciplines (football, volleyball, basketball, handball) and 2
individual disciplines (kung fu and taekwondo). The study evaluated
motivational factors affecting sport participation based on eight
fields: achievement/status, teamwork, fitness, energy release,
situational factors, skill development, friendship, and fun. Each field
was evaluated based on a survey in which every item was scored on a
three-level likert scale (very important=3, somewhat important=2, and
not important=1). A field contains several items, and an athlete's score
for that field is the \textbf{sum} of their item responses - so field
scores range well above 3. The results of the study are summarized in
the table below:

\begin{table}[H]
\centering\centering
\begin{tabular}{llrrr}
\toprule
Motivational Factor & Sport Type & Sample Mean & Sample Stdev. & Sample Size\\
\midrule
\cellcolor{gray!10}{Participation (total)} & \cellcolor{gray!10}{Team} & \cellcolor{gray!10}{75.67} & \cellcolor{gray!10}{9.38} & \cellcolor{gray!10}{203}\\
 & Individual & 80.50 & 8.40 & 62\\
\cellcolor{gray!10}{Achievement} & \cellcolor{gray!10}{Team} & \cellcolor{gray!10}{15.03} & \cellcolor{gray!10}{2.71} & \cellcolor{gray!10}{203}\\
 & Individual & 15.72 & 3.03 & 62\\
\cellcolor{gray!10}{Teamwork} & \cellcolor{gray!10}{Team} & \cellcolor{gray!10}{8.11} & \cellcolor{gray!10}{1.21} & \cellcolor{gray!10}{203}\\
\addlinespace
 & Individual & 7.68 & 1.23 & 62\\
\cellcolor{gray!10}{Energy Release} & \cellcolor{gray!10}{Team} & \cellcolor{gray!10}{11.88} & \cellcolor{gray!10}{2.05} & \cellcolor{gray!10}{203}\\
 & Individual & 12.85 & 2.07 & 62\\
\cellcolor{gray!10}{Fitness} & \cellcolor{gray!10}{Team} & \cellcolor{gray!10}{7.79} & \cellcolor{gray!10}{1.27} & \cellcolor{gray!10}{203}\\
 & Individual & 8.40 & 0.94 & 62\\
\addlinespace
\cellcolor{gray!10}{Situational Factors} & \cellcolor{gray!10}{Team} & \cellcolor{gray!10}{7.48} & \cellcolor{gray!10}{1.44} & \cellcolor{gray!10}{203}\\
 & Individual & 8.33 & 0.98 & 62\\
\cellcolor{gray!10}{Skill Development} & \cellcolor{gray!10}{Team} & \cellcolor{gray!10}{7.94} & \cellcolor{gray!10}{1.28} & \cellcolor{gray!10}{203}\\
 & Individual & 8.42 & 1.07 & 62\\
\cellcolor{gray!10}{Friendship} & \cellcolor{gray!10}{Team} & \cellcolor{gray!10}{9.57} & \cellcolor{gray!10}{1.78} & \cellcolor{gray!10}{203}\\
\addlinespace
 & Individual & 10.50 & 1.51 & 62\\
\cellcolor{gray!10}{Fun} & \cellcolor{gray!10}{Team} & \cellcolor{gray!10}{7.74} & \cellcolor{gray!10}{1.20} & \cellcolor{gray!10}{203}\\
 & Individual & 8.18 & 1.04 & 62\\
\bottomrule
\end{tabular}
\end{table}

Conduct a two independent samples \(t\)-test at the \(\alpha = 0.05\)
significance level to determine whether fitness as a motivating factor
is significantly different between team and individual sports. Assume
the populations have the SAME variance.

\(\bf (7.)\) A 1993 study published in the Canadian Journal of Zoology
investigated the characteristics of maturation and growth in two lowland
populations (Servotte and Thevenon) of European Common Frogs Rana
temporaria (Augert and Joly 1993) in Southeast France from 1986 - 1989.
A summary of the body lengths for male and female frogs from 1 - 4 years
in age in both populations are summarized below.

\begin{table}[H]
\centering\centering
\begin{tabular}{rlrrrrrr}
\toprule
Age (Years) & Population Location & Male Body Length (mm) & Stdev & Sample size & Female Body Length (mm) & Stdev & Sample Size\\
\midrule
\cellcolor{gray!10}{1} & \cellcolor{gray!10}{Servotte} & \cellcolor{gray!10}{32.5} & \cellcolor{gray!10}{5.3} & \cellcolor{gray!10}{45} & \cellcolor{gray!10}{33.4} & \cellcolor{gray!10}{5.4} & \cellcolor{gray!10}{76}\\
2 & Servotte & 55.8 & 4.7 & 11 & 59.7 & 6.2 & 15\\
\cellcolor{gray!10}{3} & \cellcolor{gray!10}{Servotte} & \cellcolor{gray!10}{61.2} & \cellcolor{gray!10}{5.3} & \cellcolor{gray!10}{35} & \cellcolor{gray!10}{65.6} & \cellcolor{gray!10}{5.2} & \cellcolor{gray!10}{27}\\
4 & Servotte & 64.7 & 7.3 & 12 & 69.6 & 0.5 & 15\\
\cellcolor{gray!10}{1} & \cellcolor{gray!10}{Thevenon} & \cellcolor{gray!10}{32.9} & \cellcolor{gray!10}{3.9} & \cellcolor{gray!10}{49} & \cellcolor{gray!10}{32.8} & \cellcolor{gray!10}{4.3} & \cellcolor{gray!10}{76}\\
\addlinespace
2 & Thevenon & 51.8 & 3.2 & 26 & 53.3 & 6.0 & 29\\
\cellcolor{gray!10}{3} & \cellcolor{gray!10}{Thevenon} & \cellcolor{gray!10}{59.9} & \cellcolor{gray!10}{5.4} & \cellcolor{gray!10}{27} & \cellcolor{gray!10}{59.9} & \cellcolor{gray!10}{6.4} & \cellcolor{gray!10}{21}\\
4 & Thevenon & 61.0 & 4.1 & 12 & 63.0 & 4.3 & 13\\
\bottomrule
\end{tabular}
\end{table}

Conduct a two independent samples \(t\)-test at the \(\alpha = 0.05\)
significance level to evaluate whether there is a statistically
significant difference in the mean body lengths of four-year-old male
frogs between the Servotte and Thévenon populations. Assume the
populations have DIFFERENT variances.

\newpage

\section*{References}\label{references}
\addcontentsline{toc}{section}{References}

\phantomsection\label{refs}
\begin{CSLReferences}{1}{0}
\bibitem[\citeproctext]{ref-augert1993plasticity}
Augert, Dominique, and Pierre Joly. 1993. {``Plasticity of Age at
Maturity Between Two Neighbouring Populations of the Common Frog (Rana
Temporaria l.).''} \emph{Canadian Journal of Zoology} 71 (1): 26--33.

\bibitem[\citeproctext]{ref-ball2000response}
Ball, John P, Kjell Danell, and Peter Sunesson. 2000. {``Response of a
Herbivore Community to Increased Food Quality and Quantity: An
Experiment with Nitrogen Fertilizer in a Boreal Forest.''} \emph{Journal
of Applied Ecology} 37 (2): 247--55.

\bibitem[\citeproctext]{ref-martins2020decrease}
Martins-Filho, Paulo R, Nicole P Damascena, Renata CM Lage, and Karyna B
Sposato. 2020. {``Decrease in Child Abuse Notifications During COVID-19
Outbreak: A Reason for Worry or Celebration?''} \emph{Journal of
Paediatrics and Child Health} 56 (12): 1980.

\bibitem[\citeproctext]{ref-moradi2020motivation}
Moradi, Jalil, Alireza Bahrami, and Amir Dana. 2020. {``Motivation for
Participation in Sports Based on Athletes in Team and Individual
Sports.''} \emph{Physical Culture and Sport. Studies and Research} 85
(1): 14--21.

\bibitem[\citeproctext]{ref-RyanBest_2024}
Ryan, Best, and Aaron Bycoffe. 2024. {``National\,: President: General
Election\,: 2024 Polls.''} \emph{FiveThirtyEight}.
\url{https://projects.fivethirtyeight.com/polls/president-general/2024/national/}.

\end{CSLReferences}

\end{document}
